\section{Related Work}
\label{sec:related}

\paragraph{Financial Multi-Agent Systems.}
LLM-based multi-agent frameworks for financial decision-making have proliferated rapidly. TradingAgents~\cite{tradingagents2025} employs a layered hierarchy of fundamental, sentiment, and technical analysts coordinated through bull-bear debate and a risk management team. FinCon~\cite{fincon2024} introduces a manager-analyst communication structure with conceptual verbal reinforcement, where self-critiquing updates investment beliefs that propagate selectively across agents. FinAgent~\cite{finagent2024} incorporates multimodal inputs but treats chart images as unaligned screenshots without structured geometric extraction. R\&D-Agent-Quant~\cite{rdagent2025} organizes factor research into a Research-Development-Feedback loop with multi-armed bandit scheduling. AlphaAgent~\cite{alphaagent2025} focuses on factor mining through AST-based originality constraints. These systems share a common limitation: their improvement mechanisms operate exclusively at the prompt level, leaving model parameters untouched by trading outcomes.

\paragraph{Self-Evolving LLM Agents.}
Beyond finance, the general agent community has developed trajectory-level parametric self-evolution methods. SE-Agent~\cite{seagent2025} applies revision, recombination, and refinement operations on agent trajectories. SEEA-R1~\cite{seea2025} combines Monte Carlo tree search with group relative policy optimization (GRPO) for iterative improvement. Self-Challenging Agents~\cite{selfchallenge2025} and Agent0~\cite{agent02026} leverage self-play and autonomous data generation. These methods demonstrate that prompt-only evolution is consistently outperformed by parametric alternatives, yet none has been adapted to the financial domain where rewards are sparse, high-variance, and regime-dependent.

\paragraph{On-Policy Self-Distillation.}
On-policy distillation has emerged as a powerful post-training paradigm. Self-Distilled Reasoner (OPSD)~\cite{opsd2026} shows that a single model can serve as both teacher and student through differential conditioning: the teacher observes privileged information (ground-truth solutions) while the student must approximate the teacher's distribution via token-level JSD matching. CRISP~\cite{crisp2026} extends this to iterative self-policy distillation. ExPO~\cite{expo2025} bootstraps from high-likelihood positive samples. Generalized Knowledge Distillation~\cite{gkd2024} establishes the on-policy foundation. Our work is the first to transfer on-policy self-distillation to financial multi-agent trajectories, using hindsight risk-adjusted PnL as privileged information and introducing Shapley credit assignment to handle the multi-agent setting.
